\documentclass[10pt,a4paper]{article}
\usepackage[margin=20mm]{geometry}
\usepackage[T1]{fontenc}
\usepackage{lmodern}
\usepackage{microtype}
\usepackage{amsmath,amssymb,amsthm,mathtools,mathrsfs}
\usepackage{graphicx}
\usepackage{xcolor}
\usepackage{booktabs}
\usepackage{enumitem}
\usepackage[hidelinks]{hyperref}

\definecolor{navy}{RGB}{21,55,94}
\definecolor{softblue}{RGB}{238,245,252}
\newtheorem{theorem}{Theorem}
\newtheorem{lemma}{Lemma}
\newtheorem{remark}{Remark}
\newcommand{\R}{\mathbb R}
\newcommand{\C}{\mathbb C}
\newcommand{\cL}{\mathcal L}
\newcommand{\cK}{\mathcal K}
\newcommand{\dd}{\,\mathrm d}
\newcommand{\supp}{\operatorname{supp}}

\title{\vspace{-1.2em}\textbf{Removing the extra smoothness unit for\\
Métivier sub-Laplacians with drift}}
\author{Candidate solution packet for arXiv:2605.02556}
\date{17 August 2026}

\begin{document}
\maketitle
\vspace{-1.2em}

\begin{center}
\fcolorbox{navy}{softblue}{\begin{minipage}{0.91\linewidth}
\textbf{Outcome.}  Let $G$ be a Métivier group of topological dimension $d$.
The multiplier theorem with drift holds under
\[
  s>d\lvert 1/p-1/2\rvert,
\]
removing the extra $+1$ in the source's condition
$s>(d+1)\lvert1/p-1/2\rvert$.  The new observation is geometric: the
positive character localizes a large propagation ball to a thin cap, whose
polynomial volume is much smaller than that of a whole outer annulus.
\end{minipage}}
\end{center}

\section*{1. The source question}

Write the stratified Lie algebra as $\mathfrak g=\mathfrak g_1\oplus
\mathfrak g_2$, with $d_i=\dim\mathfrak g_i$, $d=d_1+d_2$, and homogeneous
dimension $Q=d_1+2d_2$.  For a nontrivial positive character $\chi$, let
$\dd\mu_X=\chi\,\dd g$, $b_X=|X|/2$, and
$\mathfrak D_X=(\cL_X-b_X^2)^{1/2}$.  Garg--Singh prove the multiplier theorem
with the threshold $(d+1)|1/p-1/2|$ and ask for $d|1/p-1/2|$.

\begin{figure}[h]
\centering
\includegraphics[width=0.94\linewidth]{figures/open_question_crop.png}
\caption{The exact question on source PDF page 3.}
\end{figure}

The notation $\mathscr L_{s,\infty}^{q,-1/q,\Theta}$ below is exactly the
source's scale-invariant Sobolev class of even functions holomorphic in the
strip $|\operatorname{Im}z|<\Theta$, with boundary trace measured at
smoothness $s$.

\begin{theorem}[Optimal-form Métivier theorem]\label{thm:main}
Let $G$ be a Métivier group and $X\ne0$.  Let
$p\in[1,\infty]\setminus\{2\}$ and choose $q\in[2,\infty)$ by
\[
  \frac1q=\frac{|2/p-1|}{2},\qquad
  \Theta_{X,p}=|2/p-1|b_X.
\]
If
\[
 F\in\mathscr L_{s,\infty}^{q,-1/q,\Theta_{X,p}}
 \quad\text{for some}\quad s>d|1/p-1/2|,
\]
then $F(\mathfrak D_X)$ is bounded on $L^p(\mu_X)$ when $1<p<\infty$;
for $p=1$ it maps the local Hardy space $\mathfrak h^1(\mu_X)$ to
$L^1(\mu_X)$, and for $p=\infty$ it maps $L^\infty(\mu_X)$ to
$\mathfrak{bmo}(\mu_X)$.
\end{theorem}

\paragraph{Proof intuition.}
The source's local kernel estimates already require only $s>d/2$.  Its global
estimate uses finite propagation and then bounds a character-weighted ball by
an entire homogeneous annulus.  But $\chi(x,u)=e^{\langle a,x\rangle}$,
$|a|=2b_X$: after the exponential factor $e^{2b_Xr}$ is removed, only a thin
cap near $x=re_1$ contributes.  The cap's central fibers shrink as the square
root of the distance to the boundary.  Keeping this shrinkage removes the
extra summability loss.

\section*{2. The cap estimate}

Let $\rho$ be the Carnot--Carathéodory distance and $B_\rho(e,r)$ its ball.

\begin{lemma}[Character-weighted cap]\label{lem:cap}
For $r\ge1$ and $0\le\beta<d_2/2$,
\begin{equation}\label{eq:cap}
 \int_{B_\rho(e,r)}\frac{\chi(x,u)}{|x|^{2\beta}}\,\dd x\,\dd u
 \le C e^{2b_Xr}
 r^{(d_1+3d_2-1)/2-2\beta}.
\end{equation}
\end{lemma}

\begin{proof}
Every positive character kills the commutator layer, so
$\chi(x,u)=e^{\langle a,x\rangle}$ with $|a|=2b_X$.  Rotate the first layer so
$a=2b_Xe_1$ and write $x=te_1+z$.

We first record a fiber bound.  If $(x,u)\in B_\rho(e,r)$, choose a horizontal
path $\gamma:[0,r]\to\mathfrak g_1$ of speed at most one, with
$\gamma(0)=0$ and $\gamma(r)=x$.  Put $v=x/r$ and
$q=\gamma'-v$.  Since $\int_0^r q=0$,
\[
 \int_0^r|q|^2\dd t
 \le r-\frac{|x|^2}{r}=\frac{r^2-|x|^2}{r}.
\]
The step-two endpoint formula is
$u=\frac12\int_0^r[\gamma(t),\gamma'(t)]\dd t$.  Subtracting the straight
path and using the boundedness of the bracket gives
\[
 |u|\le C\left(r\int_0^r|q|\dd t+\Big(\int_0^r|q|\dd t\Big)^2\right)
 \le Cr\sqrt{r^2-|x|^2}.
\]
Thus, for fixed $|x|<r$, the allowed central fiber has volume at most
\begin{equation}\label{eq:fibre}
 C r^{d_2}(r^2-|x|^2)^{d_2/2}.
\end{equation}

The region $t\le r/2$ contributes at most
$Ce^{b_Xr}r^{Q-2\beta}$; for $r\ge1$ the spare exponential absorbs the
polynomial difference and makes this no larger than the right side of
\eqref{eq:cap}.  On the cap $r/2<t<r$, we have
$|x|^{-2\beta}\le Cr^{-2\beta}$.  Using \eqref{eq:fibre}, integrating first
in $z\in\R^{d_1-1}$, and then setting $y=r-t$, gives
\begin{align*}
 &Cr^{d_2-2\beta}\int_{r/2}^{r}e^{2b_Xt}
 (r^2-t^2)^{(d_1+d_2-1)/2}\dd t\\
 &\quad\le Ce^{2b_Xr}
 r^{d_2-2\beta+(d_1+d_2-1)/2}
 \int_0^\infty e^{-2b_Xy}y^{(d-1)/2}\dd y,
\end{align*}
which is exactly \eqref{eq:cap}.
\end{proof}

\section*{3. From the cap to the multiplier theorem}

The source proves, for every $0\le\beta<d_2/2$, the first-layer weighted
Plancherel estimate
\begin{equation}\label{eq:wp}
 \big\||x|^\beta\cK_{F(\sqrt\cL)}\big\|_2
 \le C\left(\int_0^\infty |F(\lambda)|^2
 \lambda^{Q-2\beta}\frac{\dd\lambda}{\lambda}\right)^{1/2}.
\end{equation}
If $\supp\widehat F\subset[-r,r]$, finite propagation puts the convolution
kernel in $B_\rho(e,r)$.  Cauchy--Schwarz, Lemma~\ref{lem:cap}, and
\eqref{eq:wp} therefore yield
\begin{equation}\label{eq:global}
 \|\chi^{1/2}\cK_{F(\sqrt\cL)}\|_1
 \le Ce^{b_Xr}(1+r)^\vartheta
 \|(1+|\cdot|)^aF\|_2,
\end{equation}
where
\begin{equation}\label{eq:parameters}
 \vartheta=\frac{d_1+3d_2-1}{4}-\beta,
 \qquad a=\frac Q2-\beta-\frac12.
\end{equation}

Now apply the conditional multiplier theorem of Martini--Ottazzi--Vallarino.
With Fourier exponent $2$, its required endpoint smoothness is any number
strictly larger than
\begin{equation}\label{eq:max}
 \max\{\sigma,\ \vartheta+1,\ a+1/2\}.
\end{equation}
Here $\sigma>d/2$ is the local kernel threshold; the source has already proved
the two required local estimates for every such $\sigma$.  Letting
$\beta\uparrow d_2/2$ in \eqref{eq:parameters},
\[
 a+\frac12\longrightarrow\frac d2,
 \qquad
 \vartheta+1\longrightarrow\frac{d+3}{4}\le\frac d2,
\]
because a Métivier group has $d_1>d_2\ge1$, hence $d\ge3$.  Thus the infimum
of \eqref{eq:max} is $d/2$.  The 2019 theorem multiplies this endpoint by
$2|1/p-1/2|$, giving precisely the hypothesis in Theorem~\ref{thm:main}.
\qed

\paragraph{Integer-derivative corollary.}
The same argument and the standard embedding used by the source replace its
Corollary 1.3 condition by
\[
 F\in H^\infty(S_{\Theta_{X,p}};N),
 \qquad N>d|1/p-1/2|.
\]

\section*{4. What is and is not resolved}

The proof fully removes the extra unit for the Métivier class studied in the
source.  Lemma~\ref{lem:cap} is purely step-two geometry and in fact holds for
every step-two stratified group.  Consequently it also gives a conditional
transfer principle: the topological threshold follows whenever the local
$d/2$ kernel estimates and the first-layer weighted Plancherel bound
\eqref{eq:wp} are available up to $\beta<d_2/2$.

The unrestricted question in the source is phrased for arbitrary two-step
groups.  In degenerate groups, the source's proof of \eqref{eq:wp} uses
precisely the Métivier nondegeneracy that is absent; the corresponding sharp
no-drift multiplier theorem is itself not known in this generality.  No claim
is made for those groups.

\paragraph{Upgrade audit.}
Six focused routes were checked: locating the loss; square-summing Fourier
time pieces; testing a radial counterexample mechanism; proving the cap
geometry; inserting its exponents into the abstract theorem; and attempting
the arbitrary two-step upgrade.  The cap route closes the Métivier case.  The
last upgrade stops at the missing weighted Plancherel estimate, not at a
summability defect.

\paragraph{Novelty boundary.}
Bounded primary-source searches through 17 August 2026 found the source,
the 2019 abstract theorem, and adjacent sharp no-drift results, but no later
paper removing the $+1$ in the drift theorem.  The argument should still be
checked by an expert in subelliptic spectral multipliers before dissemination.

\begin{thebibliography}{9}\small
\bibitem{GS}
N.~Garg and J.~Singh,
\emph{On Spectral multiplier theorem for sub-Laplacians with drift on
Métivier groups}, arXiv:2605.02556 (2026).

\bibitem{MOV}
A.~Martini, A.~Ottazzi, and M.~Vallarino,
\emph{A multiplier theorem for sub-Laplacians with drift on Lie groups},
arXiv:1705.04752; Rev. Mat. Iberoam. 35 (2019), 1501--1534.

\bibitem{N}
L.~Niedorf,
\emph{Spectral multipliers on Métivier groups}, arXiv:2412.07920.
\end{thebibliography}

\end{document}
